% Generated draft: requires mathematical and source review.
\documentclass[11pt]{article}
\usepackage[margin=25mm]{geometry}
\usepackage{fontspec}
\setmainfont{Latin Modern Roman}
\usepackage{amsmath,amssymb,mathtools,unicode-math}
\setmathfont{Latin Modern Math}
\usepackage{enumitem,needspace,xurl,hyperref}
\setlength{\parindent}{0pt}
\setlength{\parskip}{5pt}
\setlength{\emergencystretch}{3em}
\newcommand{\N}{\mathbb{N}}
\newcommand{\Z}{\mathbb{Z}}
\newcommand{\Q}{\mathbb{Q}}
\newcommand{\R}{\mathbb{R}}
\newcommand{\C}{\mathbb{C}}
\newcommand{\F}{\mathbb{F}}
\newcommand{\A}{\mathbb{A}}
\newcommand{\PP}{\mathbb{P}}
\newcommand{\E}{\mathbb{E}}
\newcommand{\eps}{\varepsilon}
\newcommand{\dd}{\,\mathrm{d}}
\begin{document}
\section*{UnsolvedMath ID: 30004723}
\textbf{Source problem code:} catalogue entry
\subsection{Definitions and mathematical statement}
A manifold is understood here to be a finite-dimensional smooth manifold. A manifold is closed if it is compact and has no boundary, and simply connected if it is path connected and has trivial fundamental group. For such a manifold \(M\), let
\[
\mathcal R_{\operatorname{Ric}>0}(M)
=
\left\{g:\ g\text{ is a smooth Riemannian metric on }M,\quad
\operatorname{Ric}_{g,p}(v,v)>0\text{ for every }p\in M\text{ and }0\ne v\in T_pM\right\}.
\]
Let \(\operatorname{Core}(M)\) stand for the assertion that \(M\) admits a core metric in the sense intended by the supplied source. The record does not define a core metric, including its geometric requirements or any boundary and normal conventions. Thus \(\operatorname{Core}(M)\) is an unresolved source-dependent predicate here, and the first question is not yet a self-contained geometric formulation.

The supplied problem asks whether both of the following assertions hold:
\begin{enumerate}
\item For every closed simply connected smooth manifold \(M\),
\[
\mathcal R_{\operatorname{Ric}>0}(M)\ne\varnothing
\quad\Longrightarrow\quad
\operatorname{Core}(M).
\]
\item For every closed simply connected smooth manifold \(M\) with \(\dim M=6\),
\[
\mathcal R_{\operatorname{Ric}>0}(M)\ne\varnothing.
\]
\end{enumerate}
The first assertion concerns existence of a core metric; it does not require the originally assumed positive-Ricci-curvature metric itself to be a core metric. No lower dimension bound is supplied for the first assertion. The smooth interpretation is the usual setting for Riemannian metrics, but the record does not explicitly specify a manifold category or address dependence on smooth structure.
\subsection{Short English statement}
Does every closed simply connected manifold admitting a positive-Ricci-curvature metric admit a core metric, and does every closed simply connected 6-manifold admit a positive-Ricci-curvature metric? These are two universal existence questions. The meaning of ``core metric'' must be recovered from the cited source before the first question is fully specified. The supplied record labels both assertions open; this status is not independently verified.
\subsection{Sources}
\begin{itemize}
\item Local registry: \texttt{data/unsolvedmath/problems.json}, supplied record with permanent UnsolvedMath ID \(30004723\), source problem code \texttt{catalogue entry}, entitled ``Core Metrics and Positive Ricci Curvature on Simply Connected Manifolds.'' Collection: Oberwolfach Reports Open Problems. The record marks the statement as exact and reports byte-for-byte retention; this draft does not independently verify that designation.
\item Supplied report citation: \emph{Analysis, Geometry and Topology of Positive Scalar Curvature Metrics} (2022), Oberwolfach Reports, pp.~1579--1649, DOI \texttt{10.4171/owr/2021/30}. \url{https://doi.org/10.4171/owr/2021/30}. The record describes Oberwolfach Reports as documenting workshops and their open-problem discussions. It gives proposed year 2021 and supplies no named proposer.
\item Supplied literature reference: ``Reiser, Journal of Topology 2024.'' \url{https://doi.org/10.1112/topo.70007}.
\item Supplied literature reference: \emph{Twisted suspensions and positive Ricci curvature}. \url{https://arxiv.org/abs/2401.06587}.
\item Literature assessment dated \texttt{2026-08-22}: the supplied curation labels the problem open and states that many new families with positive Ricci curvature and core metrics are known, while the universal core-metric assertion and the assertion for every simply connected six-manifold remain open. This is supplied dated curation derived from the OWR record, not an independently verified literature assessment. The record expressly requests independent verification of time-sensitive status and cited claims.
\end{itemize}
\end{document}
